\documentclass[aspectratio=169,11pt]{ctexbeamer}

\usetheme{Warsaw}
\setbeamertemplate{navigation symbols}{}

\usepackage{amsmath}
\usepackage{booktabs}
\usepackage{tabularx}
\usepackage{fancyvrb}
\usepackage{xcolor}
\usepackage{tikz}
\usetikzlibrary{arrows.meta,positioning}

\definecolor{codebg}{RGB}{248,248,248}
\definecolor{deepblue}{RGB}{37,99,235}
\definecolor{darkgreen}{RGB}{22,101,52}
\definecolor{warmorange}{RGB}{217,119,6}
\setbeamercolor{block body}{bg=codebg,fg=black}
\setbeamercolor{block title}{bg=deepblue,fg=white}

\title{三维绘图与可视化基本概念}
\subtitle{围绕 Plotly/Dash 单位立方体与单位球示例}
\author{crazyfish}
\date{\today}

\begin{document}

\begin{frame}
  \titlepage
\end{frame}

\begin{frame}{本讲目标}
  通过一个最小三维场景理解三维绘图的核心组成。

  \begin{itemize}
    \item 物体：单位立方体与单位球。
    \item 空间：坐标轴、尺度、中心与方向。
    \item 观察：相机位置、视角与投影。
    \item 渲染：颜色、光源、漫反射、镜面反射与粗糙度。
    \item 交互：滑杆、鼠标旋转、实时更新。
  \end{itemize}
\end{frame}

\begin{frame}[fragile]{运行课堂示例}
  本讲围绕当前目录中的 Dash + Plotly 示例展开。

  \begin{block}{启动命令}
\begin{Verbatim}[fontsize=\small]
conda activate ai4math-vis
python basic_3d/plotly_unit_cube_demo.py
\end{Verbatim}
  \end{block}

  \begin{block}{浏览器地址}
\begin{Verbatim}[fontsize=\small]
http://127.0.0.1:8051
\end{Verbatim}
  \end{block}

  \begin{itemize}
    \item \texttt{Object} 可切换单位立方体与单位球。
    \item \texttt{Lighting} 控制光源位置与材质响应。
    \item 鼠标拖动三维图会同步更新相机位置滑杆。
  \end{itemize}
\end{frame}

\section{三维数据}

\begin{frame}{三维绘图首先是空间数据}
  三维图不是“二维图变立体”，而是数据中真的包含空间结构。

  \begin{table}
    \centering
    \begin{tabularx}{0.9\linewidth}{lllX}
      \toprule
      对象 & 数据 & Plotly 图元 & 典型用途 \\
      \midrule
      点 & \((x,y,z)\) & \texttt{Scatter3d} & 采样点、粒子、传感器 \\
      线 & 点序列 & \texttt{Scatter3d} & 轨迹、边框、流线 \\
      面 & 网格或三角片 & \texttt{Surface/Mesh3d} & 几何边界、地形、等值面 \\
      体 & 三维数组 & \texttt{Volume/Isosurface} & 医学、流体、标量场 \\
      \bottomrule
    \end{tabularx}
  \end{table}
\end{frame}

\begin{frame}{坐标系：三维可视化的参照物}
  三维对象需要放在明确的坐标系中。

  \begin{itemize}
    \item 坐标轴给出方向：\(x\)、\(y\)、\(z\)。
    \item 刻度给出尺度：一个单位到底有多大。
    \item 中心位置决定光源、相机和旋转的参照。
    \item 当前示例中：
      \begin{itemize}
        \item 单位立方体：\([0,1]\times[0,1]\times[0,1]\)。
        \item 单位球：球心 \((0,0,0)\)，半径 \(1\)。
      \end{itemize}
  \end{itemize}
\end{frame}

\begin{frame}{单位立方体：最小的三维几何模型}
  单位立方体适合解释三维几何的基本元素。

  \begin{columns}
    \begin{column}{0.48\linewidth}
      \begin{itemize}
        \item 8 个顶点。
        \item 12 条边。
        \item 6 个面。
        \item 每个面有固定法向。
      \end{itemize}
    \end{column}
    \begin{column}{0.48\linewidth}
      \begin{block}{六个面的法向}
        \[
        \pm\mathbf e_x,\quad
        \pm\mathbf e_y,\quad
        \pm\mathbf e_z
        \]
      \end{block}
      面法向是判断“面是否朝向光源”的关键。
    \end{column}
  \end{columns}
\end{frame}

\begin{frame}{单位球：连续曲面的代表}
  单位球适合解释曲面、法向连续变化和高光。

  \begin{block}{单位球方程}
    \[
    x^2+y^2+z^2=1
    \]
  \end{block}

  \begin{itemize}
    \item 球面上每一点都有自己的法向。
    \item 对单位球，点 \(\mathbf p=(x,y,z)\) 的法向就是 \(\mathbf n=\mathbf p\)。
    \item 因为法向连续变化，漫反射和镜面反射在球面上更容易观察。
  \end{itemize}
\end{frame}

\section{绘制管线}

\begin{frame}{三维绘图管线}
  从数学对象到屏幕图像，一般经历如下步骤。

  \begin{center}
    \begin{tikzpicture}[
      node distance=0.75cm,
      box/.style={draw, rounded corners=2pt, align=center, minimum width=1.75cm, minimum height=0.72cm},
      arr/.style={-{Stealth[length=2.2mm]}, thick}
    ]
      \node[box] (data) {数据\\坐标/网格};
      \node[box, right=of data] (geom) {几何\\点/线/面};
      \node[box, right=of geom] (view) {相机\\视角/投影};
      \node[box, right=of view] (shade) {光照\\颜色/材质};
      \node[box, right=of shade] (screen) {屏幕\\交互显示};
      \draw[arr] (data) -- (geom);
      \draw[arr] (geom) -- (view);
      \draw[arr] (view) -- (shade);
      \draw[arr] (shade) -- (screen);
    \end{tikzpicture}
  \end{center}

  \begin{itemize}
    \item Plotly 负责把图元渲染到浏览器中的 WebGL 画布。
    \item Dash 负责把滑杆、下拉框等控件连接到 Python 回调。
  \end{itemize}
\end{frame}

\begin{frame}[fragile]{Plotly 图元：\texttt{Surface} 与 \texttt{Scatter3d}}
  \footnotesize
  当前示例主要使用两个三维图元。

  \begin{block}{球面或细分面}
\begin{Verbatim}[fontsize=\scriptsize]
go.Surface(
    x=x_values,
    y=y_values,
    z=z_values,
    surfacecolor=intensities,
)
\end{Verbatim}
  \end{block}

  \begin{block}{光源点与虚线}
\begin{Verbatim}[fontsize=\scriptsize]
go.Scatter3d(
    x=[light_x, target_x],
    y=[light_y, target_y],
    z=[light_z, target_z],
    mode="markers+lines+text",
)
\end{Verbatim}
  \end{block}
\end{frame}

%% \begin{frame}{为什么不用 Plotly 内部 \texttt{lightposition}}
%%   本例采用显式光照计算，而不是完全依赖 Plotly 内部光照。

%%   \begin{itemize}
%%     \item Plotly 的三维光照会经过内部矩阵变换。
%%     \item 可见的“光源点”和内部 shader 光源方向不一定直观一致。
%%     \item 教学时更需要“光源在哪边，哪边更亮”的稳定对应关系。
%%     \item 因此本例直接计算每个面或每个采样点的亮度，再把亮度映射为颜色。
%%   \end{itemize}
%% \end{frame}

\section{相机}

\begin{frame}{相机：决定我们从哪里看}
  三维场景最终仍显示在二维屏幕上，相机定义观察方式。

  \begin{itemize}
    \item \texttt{camera.eye}：相机所在位置。
    \item \texttt{camera.center}：观察中心偏移。
    \item \texttt{camera.up}：画面中“上方”的方向。
    \item 鼠标旋转三维图，本质上是在改变相机视角。
  \end{itemize}

  \begin{block}{当前示例}
    鼠标调整相机后，\texttt{Camera X/Y/Z} 滑杆会同步更新。
  \end{block}
\end{frame}

\begin{frame}[fragile]{相机参数示意}
  Plotly 中的三维相机通常写在 \texttt{layout.scene.camera} 中。

  \begin{block}{示例结构}
\begin{Verbatim}[fontsize=\small]
scene={
    "camera": {
        "eye": {"x": camera_x, "y": camera_y, "z": camera_z},
        "center": {"x": 0, "y": 0, "z": 0},
    }
}
\end{Verbatim}
  \end{block}

  \begin{itemize}
    \item 相机离物体越远，物体看起来越小。
    \item 相机绕物体移动，看到的亮面和高光也会变化。
  \end{itemize}
\end{frame}

\section{光照与材质}

\begin{frame}{光照模型的基本思想}
  当前示例使用简化的局部光照模型。

  \begin{block}{核心公式}
    \[
    I =
    A
    + D\max(\mathbf n\cdot \mathbf l,0)
    + S\,H
    + F\,E
    \]
  \end{block}

  \begin{itemize}
    \item \(\mathbf n\)：表面法向。
    \item \(\mathbf l\)：从表面指向光源的单位方向。
    \item \(H\)：镜面高光项。
    \item \(E\)：边缘视角反光项。
  \end{itemize}
\end{frame}

\begin{frame}{Ambient：环境光}
  \texttt{Ambient 环境光} 不是光源强度，而是基础亮度。

  \begin{itemize}
    \item 不依赖光源方向。
    \item 所有面都会获得一部分亮度。
    \item 值越大，阴影越浅，对比越弱。
    \item 值越小，背光面越暗，光照方向更明显。
  \end{itemize}

  \begin{block}{教学观察}
    把 \texttt{Ambient} 降低，再移动光源，光照方向最明显。
  \end{block}
\end{frame}

\begin{frame}{Diffuse：漫反射}
  \texttt{Diffuse 漫反射} 是最直观的受光效果。

  \begin{block}{Lambert 漫反射}
    \[
    I_D = D\max(\mathbf n\cdot \mathbf l,0)
    \]
  \end{block}

  \begin{itemize}
    \item 面朝向光源：\(\mathbf n\cdot\mathbf l\) 大，亮。
    \item 面背向光源：\(\mathbf n\cdot\mathbf l<0\)，不接受漫反射。
    \item 单位立方体中，移动光源可以直接观察哪个面变亮。
  \end{itemize}
\end{frame}

\begin{frame}{Specular 与 Roughness：高光和粗糙度}
  \texttt{Specular 镜面反射} 控制高光强度，\texttt{Roughness 粗糙度} 控制高光范围。

  \begin{itemize}
    \item \texttt{Specular} 越大，高光越亮。
    \item \texttt{Roughness} 越小，高光越集中。
    \item \texttt{Roughness} 越大，高光越宽、越钝。
    \item 球面比立方体更适合观察连续高光。
  \end{itemize}

  \begin{block}{课堂操作}
    切换到单位球，把 \texttt{Specular} 从 0 拉到 2，再调节 \texttt{Roughness}。
  \end{block}
\end{frame}

\begin{frame}{Fresnel：斜视角反光}
  \texttt{Fresnel 菲涅尔} 描述反射强度随观察角度变化。

  \begin{itemize}
    \item 正对表面时，边缘反光较弱。
    \item 接近掠射角观察时，边缘反光更强。
    \item 在球面边缘区域更容易观察。
    \item 本例中的 Fresnel 是近似，不是完整物理渲染。
  \end{itemize}
\end{frame}

\section{交互}

\begin{frame}{为什么要交互}
  三维概念很难只靠静态图说明。

  \begin{itemize}
    \item 旋转相机：理解三维结构。
    \item 移动光源：理解法向和受光方向。
    \item 调材质参数：理解渲染风格。
    \item 切换对象：比较平面几何与曲面几何。
  \end{itemize}

  \begin{block}{本例的关键交互}
    所有滑杆使用拖动时更新，参数变化可以连续观察。
  \end{block}
\end{frame}

%% \begin{frame}[fragile]{Dash 回调：控件驱动图像更新}
%%   Dash 的核心思想是：输入控件变化，回调函数重新生成图形。

%%   \begin{block}{概念结构}
%% \begin{Verbatim}[fontsize=\small]
%% @app.callback(
%%     Output("cube-graph", "figure"),
%%     Input("object-shape", "value"),
%%     Input("light-x", "value"),
%%     Input("ambient", "value"),
%%     ...
%% )
%% def update_cube(...):
%%     controls = normalize_controls(...)
%%     return make_figure_from_controls(controls)
%% \end{Verbatim}
%%   \end{block}
%% \end{frame}

\begin{frame}{演示顺序建议}
  \begin{enumerate}
    \item 先只看单位立方体，打开坐标轴，关闭边框。
    \item 移动光源，观察哪一个面变亮。
    \item 降低 \texttt{Ambient}，增强明暗对比。
    \item 切换到单位球，观察连续明暗变化。
    \item 调节 \texttt{Specular} 和 \texttt{Roughness}，观察高光。
    \item 鼠标旋转相机，观察相机位置对高光的影响。
  \end{enumerate}
\end{frame}

\begin{frame}{从这个例子推广到模拟结果可视化}
  计算模拟结果通常只是把“几何对象”换成更复杂的数据。

  \begin{itemize}
    \item 地形：二维高度场 \(\rightarrow\) 曲面。
    \item 流体：速度向量 \(\rightarrow\) 箭头、流线。
    \item 标量场：温度/密度 \(\rightarrow\) 切片、等值面、体渲染。
    \item 医学图像：三维体素 \(\rightarrow\) 切片与体渲染。
  \end{itemize}

  \begin{block}{不变的核心概念}
    坐标、网格、相机、颜色映射、光照、交互。
  \end{block}
\end{frame}

\begin{frame}{小结}
  三维绘图的基本概念可以用一个很小的场景讲清楚。

  \begin{itemize}
    \item 几何决定物体形状。
    \item 坐标轴决定空间参照。
    \item 相机决定观察方式。
    \item 光源和材质决定明暗与高光。
    \item 交互让参数和视觉结果建立直接联系。
  \end{itemize}

  \begin{block}{本讲示例文件}
    \texttt{basic\_3d/plotly\_unit\_cube\_demo.py}
  \end{block}
\end{frame}

\end{document}
